% preamble-exam-zh.tex — 极简讲解页共享 preamble
% 目标：复刻「题目独立、提示轻框、路线箭头、主体双栏、答案轻收束」的讲义风格。

\documentclass{exam-zh}

% ---------- 基础配置 ----------
\examsetup{
  page/size = a4paper,
  page/show-head = false,
  page/show-foot = false,
  sealline/show = false,
  font = none,
  math-font = none,
}

% ---------- 包 ----------
\usepackage{geometry}
\geometry{top=10mm,bottom=14mm,left=15mm,right=13mm}
\AtBeginDocument{\newgeometry{top=10mm,bottom=14mm,left=15mm,right=13mm}}

\usepackage{xcolor}
\usepackage{needspace}
\usepackage{enumitem}
\usepackage{array}
\usepackage{tabularx}
\usepackage{tcolorbox}
\tcbuselibrary{skins,breakable}
\usepackage{paracol}

% ---------- 打印/屏幕颜色切换 ----------
% 用法：\documentclass[monochrome]{exam-zh} 可降级为灰度。
% 注意：不要使用 \if@xxx 放在 makeatother 外部；这里改成普通条件名。
\newif\ifeduprintcolor
\eduprintcolortrue
\makeatletter
\@ifclasswith{exam-zh}{monochrome}{\eduprintcolorfalse}{}
\makeatother

% ---------- 语义颜色 ----------
\ifeduprintcolor
  \definecolor{edu-blue}{RGB}{31,62,126}
  \definecolor{edu-blue-soft}{RGB}{232,238,250}
  \definecolor{edu-blue-bg}{RGB}{248,250,254}
  \definecolor{edu-ink}{RGB}{30,35,45}
  \definecolor{edu-muted}{RGB}{118,124,136}
  \definecolor{edu-line}{RGB}{209,218,235}
  \definecolor{edu-soft-bg}{RGB}{250,251,253}
  \definecolor{edu-red}{RGB}{168,58,58}
  \definecolor{edu-red-bg}{RGB}{255,247,245}
\else
  \definecolor{edu-blue}{gray}{0.15}
  \definecolor{edu-blue-soft}{gray}{0.90}
  \definecolor{edu-blue-bg}{gray}{0.98}
  \definecolor{edu-ink}{gray}{0.10}
  \definecolor{edu-muted}{gray}{0.45}
  \definecolor{edu-line}{gray}{0.82}
  \definecolor{edu-soft-bg}{gray}{0.97}
  \definecolor{edu-red}{gray}{0.28}
  \definecolor{edu-red-bg}{gray}{0.97}
\fi

% ---------- 全局排版 ----------
\setlength{\parindent}{0pt}
\setlength{\parskip}{0pt}
\linespread{1.16}
\renewcommand{\arraystretch}{1.15}

% ctex 通常提供 \kaishu；这里用它做教师旁白感。
\newcommand{\edukai}{\kaishu}

% ---------- 顶部元信息 ----------
\newcommand{\eduTopBar}[2]{%
  \noindent
  \begin{tabularx}{\linewidth}{@{}Xr@{}}
    {\color{edu-blue}\bfseries #1} &
    {\color{edu-blue}\bfseries #2}
  \end{tabularx}
  \vspace{8mm}
}

% ---------- 轻标题体系 ----------
% 只有真正的大结构才用它；不要给提示/路线滥用标题。
\newcommand{\eduSectionTitle}[1]{%
  \needspace{5\baselineskip}%
  \vspace{2mm}
  {\color{edu-blue}\bfseries\large #1}\par
  \vspace{3mm}
}

\newcommand{\eduSubTitle}[1]{%
  \needspace{4\baselineskip}%
  \vspace{1mm}
  {\color{edu-blue}\bfseries #1}\par
  \vspace{1mm}
}

% ---------- 小问题干：复用 exam (1)(2)(3) 格式 ----------
\newenvironment{eduSubQuestionItem}[1]{%
  \needspace{4\baselineskip}%
  \vspace{2mm}
  \par\noindent
  \hangindent=8mm\hangafter=1
  {\color{edu-blue}\bfseries #1}\enspace
  \begingroup
  \color{edu-ink}
}{%
  \endgroup
  \par
  \vspace{3mm}
}

% ---------- 辅助信息条（解题流程/关键想法/思考）楷体 ----------
\newenvironment{eduAuxItem}[1]{%
  \needspace{4\baselineskip}%
  \vspace{2mm}
  \par\noindent
  \hangindent=8mm\hangafter=1
  {\color{edu-blue}\edukai $\triangleright$~#1}\enspace
  \begingroup
  \color{edu-ink}
  \edukai
}{%
  \endgroup
  \par
  \vspace{4mm}
}

\newcommand{\eduColumnTitle}[2]{%
  {\color{edu-blue}\bfseries #1\hspace{1.5mm}#2}\par
  \vspace{2.5mm}
}

% ---------- 图标：不用外部 icon 字体，保证可移植 ----------
\newcommand{\eduIconBulb}{\(\circ\)}
\newcommand{\eduIconBook}{\(\square\)}
\newcommand{\eduIconCheck}{\(\checkmark\)}
\newcommand{\eduIconPin}{\(\diamond\)}
\newcommand{\eduIconNote}{\(\triangleright\)}

% ---------- 题目区 ----------
% 题目区是页面入口：弱标签 + 一条长蓝线 + 大题干。
\newenvironment{eduProblemBlock}[1][题目]{%
  \needspace{8\baselineskip}
  {\small\color{edu-muted}#1}\par
  \vspace{1.5mm}
  {\color{edu-blue}\hrule height 0.55pt}
  \vspace{4.5mm}
  \begingroup
  \color{edu-ink}
  \large
  \setlength{\parskip}{2.2mm}
  \linespread{1.20}\selectfont
}{%
  \par
  \endgroup
  \vspace{6mm}
}

\newenvironment{eduSubQuestions}{%
  \begin{enumerate}[
    label=(\arabic*),
    leftmargin=8mm,
    itemsep=2.0mm,
    topsep=2.0mm,
    parsep=0pt
  ]
}{%
  \end{enumerate}
}

% ---------- 读题提示：轻框，不做同级标题 ----------
\newtcolorbox{eduReadTipBox}{
  enhanced,
  breakable,
  colback=white,
  colframe=edu-blue!45,
  boxrule=0.45pt,
  arc=1.6mm,
  left=10mm,
  right=4mm,
  top=5mm,
  bottom=5mm,
  before skip=1mm,
  after skip=7mm,
  fontupper=\edukai\normalsize,
  colupper=edu-ink,
  overlay={
    \node[anchor=north west, text=edu-blue] at ([xshift=3.2mm,yshift=-3.2mm]frame.north west)
    {\small\eduIconNote};
  }
}

\newenvironment{eduTipList}{%
  \begin{itemize}[
    label=\textbullet,
    leftmargin=5mm,
    itemsep=1.2mm,
    topsep=0pt,
    parsep=0pt
  ]
}{%
  \end{itemize}
}

% ---------- 解题路线：虚线框 + 文字箭头 ----------
\newenvironment{eduRouteFlow}{%
  \needspace{4\baselineskip}%
  \vspace{1mm}
  \begin{center}
  \normalsize
}{%
  \end{center}
  \vspace{7mm}
}
\newcommand{\eduRouteStep}[1]{%
  {\color{edu-ink}#1}%
}
\newcommand{\eduRouteArrow}{%
  \;$\boldsymbol{\to}$\;%
}

% ---------- 主体讲解：使用 paracol 实现跨页自适应双栏 ----------
\newenvironment{eduDualExplain}{%
  \needspace{6\baselineskip}% 防止标题孤行
  \columnratio{0.25, 0.75}%
  \setlength{\columnsep}{7mm}%
  \columnseprule=0.45pt%
  \colseprulecolor{edu-blue}%
  \begin{paracol}{2}% 开启双栏
}{%
  \end{paracol}% 结束双栏
  \vspace{5mm}
}

\newenvironment{eduExplainLeft}{%
  \normalsize\color{edu-ink}%
}{%
  \switchcolumn
}

\newcommand{\eduColumnDivider}{}

\newenvironment{eduExplainRight}{%
  \normalsize\color{edu-ink}%
}{%
}

\newenvironment{eduNumberList}{%
  \begin{enumerate}[
    label=\arabic*.,
    leftmargin=6mm,
    itemsep=4.8mm,
    topsep=0pt,
    parsep=0pt
  ]
}{%
  \end{enumerate}
}

\newenvironment{eduBulletList}{%
  \begin{itemize}[
    label=\textbullet,
    leftmargin=5mm,
    itemsep=1.2mm,
    topsep=0pt,
    parsep=0pt
  ]
}{%
  \end{itemize}
}

% ---------- 底部双栏：答案 + 方法提醒 ----------
\newenvironment{eduBottomGrid}{%
  \vspace{1mm}
  \par\noindent
}{%
  \par\vspace{1mm}
}

\newenvironment{eduSummaryLeft}{%
  \begin{minipage}[t]{0.30\linewidth}
}{%
  \end{minipage}
}

\newcommand{\eduSummaryDivider}{%
  \hspace{4mm}{\color{edu-line}\vrule width 0.35pt}\hspace{7mm}%
}

\newenvironment{eduSummaryRight}{%
  \begin{minipage}[t]{0.58\linewidth}
}{%
  \end{minipage}
}

% ---------- 答案/方法提醒：轻量收束 ----------
\newtcolorbox{eduSummaryBlock}[2]{
  enhanced,
  breakable,
  colback=white,
  colframe=white,
  boxrule=0pt,
  sharp corners,
  left=0mm,
  right=0mm,
  top=1mm,
  bottom=1mm,
  before skip=1mm,
  after skip=1mm,
  title={\color{edu-blue}\bfseries #1\hspace{1.5mm}#2},
  coltitle=edu-blue,
  attach title to upper={\par\vspace{2mm}},
  fontupper=\normalsize
}

\newenvironment{eduPlainList}{%
  \begin{enumerate}[
    label=(\arabic*),
    leftmargin=7mm,
    itemsep=1.2mm,
    topsep=0pt,
    parsep=0pt
  ]
}{%
  \end{enumerate}
}

% ---------- 兼容旧 block 的轻量盒子 ----------
\newtcolorbox{eduSoftBox}{
  enhanced,
  breakable,
  colback=edu-soft-bg,
  colframe=edu-line,
  boxrule=0.35pt,
  arc=1mm,
  left=3mm,
  right=3mm,
  top=2mm,
  bottom=2mm,
  before skip=1mm,
  after skip=2mm,
  fontupper=\small
}

\newtcolorbox{eduMistakeBox}{
  enhanced,
  breakable,
  colback=edu-red-bg,
  colframe=edu-red!40,
  boxrule=0.35pt,
  arc=1mm,
  left=3mm,
  right=3mm,
  top=2.5mm,
  bottom=2.5mm,
  before skip=1mm,
  after skip=2mm,
  fontupper=\normalsize
}

\newtcolorbox{eduLegacySolution}{
  enhanced,
  breakable,
  colback=edu-soft-bg,
  colframe=edu-line,
  boxrule=0.35pt,
  arc=1mm,
  left=3mm,
  right=3mm,
  top=2mm,
  bottom=2mm,
  before skip=-2mm,
  after skip=4mm,
  fontupper=\small
}

\newtcolorbox{eduTeacherNote}{
  enhanced,
  breakable,
  colback=edu-soft-bg,
  colframe=edu-line,
  boxrule=0pt,
  borderline west={1.5pt}{0pt}{edu-muted},
  left=2.5mm,
  right=2.5mm,
  top=2mm,
  bottom=2mm,
  before skip=1mm,
  after skip=2mm,
  fontupper=\footnotesize,
  colupper=edu-muted
}

% ---------- 答题区 ----------
\newcommand{\answerarea}[1][40mm]{%
  \vspace{3mm}%
  \begin{tcolorbox}[
    enhanced,
    breakable,
    colback=white,
    colframe=edu-line,
    boxrule=0.55pt,
    arc=0.8mm,
    height=#1,
    valign=top,
  ]
  \end{tcolorbox}%
}


\begin{document}

\eduSectionTitle{方法总论：共轭展开与公式选择框架}





\begin{eduAuxItem}{关键想法}
对复数和、差模平方做共轭展开：
\begin{align*}
  |z_1 + z_2|^2 &= (z_1 + z_2)(\bar{z}_1 + \bar{z}_2) = |z_1|^2 + |z_2|^2 + (z_1\bar{z}_2 + \bar{z}_1z_2) \\
  |z_1 - z_2|^2 &= (z_1 - z_2)(\bar{z}_1 - \bar{z}_2) = |z_1|^2 + |z_2|^2 - (z_1\bar{z}_2 + \bar{z}_1z_2)
\end{align*}
两式相加即得\textbf{平行四边形恒等式}（消去交叉项）：
\begin{equation*}
  |z_1 + z_2|^2 + |z_1 - z_2|^2 = 2(|z_1|^2 + |z_2|^2)
\end{equation*}
两式相减可得\textbf{交叉项表达式}（消去模长和）：
\begin{equation*}
  |z_1 + z_2|^2 - |z_1 - z_2|^2 = 2(z_1\bar{z}_2 + \bar{z}_1z_2)
\end{equation*}
利用复数商的关系 $z_1/z_2$，我们可以将交叉项改写为（记 $w = \dfrac{z_1}{z_2}$）：
\begin{equation*}
  z_1\bar{z}_2 + \bar{z}_1z_2 = |z_2|^2 \cdot \frac{z_1}{z_2} + |z_1|^2 \cdot \frac{z_2}{z_1} = |z_2|^2 \left( w + \bar{w} \right) = 2|z_2|^2 \mathrm{Re}(w)
\end{equation*}\end{eduAuxItem}



\eduSectionTitle{例题 1 — 复数知三推一}

\begin{eduProblemBlock}[例题 1]
已知复数 $z_1, z_2$ 满足 $|z_1| = \sqrt{5}$，$|z_1 - z_2| = \sqrt{2}$，$|z_1 + z_2| = 3\sqrt{2}$，求 $|z_2|$ 及 $\dfrac{z_1}{z_2}$。
\end{eduProblemBlock}









\begin{eduSubQuestionItem}{第一步}
应用平行四边形公式求得 $|z_2|$ 的值。\end{eduSubQuestionItem}

\begin{eduDualExplain}
  \begin{eduExplainLeft}
    \eduColumnTitle{\eduIconBulb}{思考引导}
    \begin{eduNumberList}
      \item 我们要连结 $|z_1|$, $|z_2|$, $|z_1-z_2|$, $|z_1+z_2|$ 这四个模长。      \item 写出平行四边形恒等式，将已知的三项代入即可。    \end{eduNumberList}
  \end{eduExplainLeft}
  \eduColumnDivider
  \begin{eduExplainRight}
    \eduColumnTitle{\eduIconBook}{规范讲解}
    \begin{eduNumberList}
      \item 平行四边形恒等式为：$|z_1 - z_2|^2 + |z_1 + z_2|^2 = 2(|z_1|^2 + |z_2|^2)$。      \item 代入已知数据：$(\sqrt{2})^2 + (3\sqrt{2})^2 = 2((\sqrt{5})^2 + |z_2|^2)$。      \item 计算化简：$2 + 18 = 2(5 + |z_2|^2) \implies 20 = 10 + 2|z_2|^2$。      \item 得出 $|z_2|^2 = 5$。因为模长非负，故 $|z_2| = \sqrt{5}$。    \end{eduNumberList}

  \end{eduExplainRight}
\end{eduDualExplain}




\begin{eduSubQuestionItem}{第二步}
通过展开公式或两式相减求出 $z_1/z_2$。\end{eduSubQuestionItem}

\begin{eduDualExplain}
  \begin{eduExplainLeft}
    \eduColumnTitle{\eduIconBulb}{思考引导}
    \begin{eduNumberList}
      \item 设 $w = z_1/z_2$，我们知道 $|w| = |z_1|/|z_2| = \sqrt{5}/\sqrt{5} = 1$。      \item 要求出 $w$，我们需要先求出它的实部 $\mathrm{Re}(w) = \cos\theta$。      \item 利用模平方之差公式展开：$|z_1+z_2|^2 - |z_1-z_2|^2 = 4\mathrm{Re}(z_1\bar{z}_2) = 4|z_2|^2\mathrm{Re}(w)$。    \end{eduNumberList}
  \end{eduExplainLeft}
  \eduColumnDivider
  \begin{eduExplainRight}
    \eduColumnTitle{\eduIconBook}{规范讲解}
    \begin{eduNumberList}
      \item 由模平方之差公式得：$|z_1+z_2|^2 - |z_1-z_2|^2 = 2(z_1\bar{z}_2 + \bar{z}_1z_2)$。      \item 因为 $z_1\bar{z}_2 + \bar{z}_1z_2 = 2|z_2|^2 \mathrm{Re}(w)$（其中 $w = z_1/z_2$）。      \item 代入数据：$18 - 2 = 4|z_2|^2 \mathrm{Re}(w) \implies 16 = 4 \times 5 \mathrm{Re}(w)$。      \item 得到 $\mathrm{Re}(w) = \dfrac{16}{20} = \dfrac{4}{5}$。      \item 由于 $|w| = \dfrac{|z_1|}{|z_2|} = 1$。      \item 根据 $|w|^2 = [\mathrm{Re}(w)]^2 + [\mathrm{Im}(w)]^2 = 1$，得 $[\mathrm{Im}(w)]^2 = 1 - \left(\dfrac{4}{5}\right)^2 = \dfrac{9}{25}$。      \item 因为 $\mathrm{Im}(w) = \pm\dfrac{3}{5}$，所以 $\dfrac{z_1}{z_2} = \dfrac{4 \pm 3i}{5}$。    \end{eduNumberList}

    \vspace{1mm}
    \eduSubTitle{结论}
    \begin{eduBulletList}
      \item 两步环环相扣，最终答案为：$|z_2| = \sqrt{5}$，$\dfrac{z_1}{z_2} = \dfrac{4 \pm 3i}{5}$。    \end{eduBulletList}
  \end{eduExplainRight}
\end{eduDualExplain}




\begin{eduBottomGrid}
  \begin{eduSummaryLeft}
    \begin{eduSummaryBlock}{\eduIconCheck}{避坑指南}
    \begin{eduPlainList}
      \item **不要硬设坐标**：设 $z_1=a+bi, z_2=c+di$ 展开会得到含4个未知数的3个二次方程，极难求解！      \item **别忘比值双解**：由平方和为1开根号时，记得加 $\pm$，即 $\mathrm{Im}(z_1/z_2) = \pm 3/5$。      \item **分清角和与角差**：商 $z_1/z_2$ 的辐角是角差 $\alpha - \beta$，而 $z_1\bar{z}_2$ 也是角差；共轭相乘是角和。    \end{eduPlainList}
    \end{eduSummaryBlock}
  \end{eduSummaryLeft}
  \eduSummaryDivider
  \begin{eduSummaryRight}
    \begin{eduSummaryBlock}{\eduIconPin}{万能公式模板}
    \begin{eduBulletList}
      \item 平行四边形恒等式：$|z_1-z_2|^2+|z_1+z_2|^2 = 2(|z_1|^2+|z_2|^2)$      \item 模平方差与交叉项：$|z_1+z_2|^2-|z_1-z_2|^2 = 4\mathrm{Re}(z_1\bar{z}_2)$      \item 复商与交叉项关系：$z_1\bar{z}_2+\bar{z}_1z_2 = 2|z_2|^2\mathrm{Re}(z_1/z_2)$      \item 极坐标关联：$\mathrm{Re}(z_1/z_2) = \dfrac{|z_1|}{|z_2|}\cos\theta$    \end{eduBulletList}
    \end{eduSummaryBlock}
  \end{eduSummaryRight}
\end{eduBottomGrid}




\end{document}
